\section{Positive heights and the critical arithmetic transition}
\label{sec:v60-heights-contact}
The same perturbed roof now admits positive boundary estimates and an
explicit density profile in the moving-peak regime. The source, coded
states and stable endpoint difference are those already used for
Theorem~\ref{thm:v59-markov-contact}; no perturbation is reset to zero
before the density is evaluated.

\subsection{Three original boundary sources, uniformly in the parameter}
Keep $U_s,V_s$ and the initial-age strips of
Theorem~\ref{thm:v59-markov-boundary-height}. Let
$p_{m,\epsilon,s}^{\rm vert},p_{m,\epsilon,s}^{\rm stab}$ and
$p_{m,\epsilon,s}^{\rm age}$ denote their unnormalized positive
$S_m\tau_\epsilon$ densities. Put
$L_m=\min\{\lfloor m/2\rfloor,\lfloor\sqrt m\rfloor\}$.

\begin{theorem}[Perturbed positive essential heights]
\label{thm:v60-perturbed-height}
For $m\ge2$, $0<s<1/8$, and any fixed $0<\epsilon_0<1/4$,
\begin{equation}\label{eq:v60-perturbed-height}
 \sup_{|\epsilon|\le\epsilon_0}\sqrt m
 \left(\|p_{m,\epsilon,s}^{\rm vert}\|_\infty
       +\|p_{m,\epsilon,s}^{\rm stab}\|_\infty
       +\|p_{m,\epsilon,s}^{\rm age}\|_\infty\right)
                      \le C(s+\rho_0^{L_m}).
\end{equation}
The vertical and terminal stable terms separately satisfy the stronger
bound with $\rho_0^{\lfloor m/2\rfloor}$. The same domination applies
to bounded insertions supported on the union, with their supremum norm.
Consequently the conditional probability of this union, given the
\emph{same} roof value $t$ with
$|t-m\bar\tau_\epsilon|\le M\sqrt m$, is at most
$C_M(s+\rho_0^{L_m})$ for almost every such $t$ and all sufficiently
large $m$, uniformly in $\epsilon$.
\end{theorem}
\begin{proof}
The coarea derivative is $1-a_w\ge1-\rho_0$, and every contributing
word satisfies $|A_m+\epsilon B_m-(t-3m)|\le1$.
For a fixed prefix of depth $L=\lfloor m/2\rfloor$, subtract its
exact two rewards. The remaining strip probability is at most
$C/\sqrt m$, by \eqref{eq:v60-strip-upper}, uniformly in the prefix,
its terminal state and the shifted roof value. The outer forward
cylinder cover of $U_s$ has stationary mass at most $C(s+\rho_0^L)$.
Apply the first criterion of Proposition~\ref{prop:v60-coarea-height}.
This uses the residual local constraint after fixing the prefix.

For $V_s$ cover its terminal stable intervals by suffix cylinders of
the same depth. Their total stationary mass is at most
$C(s+\rho_0^L)$; terminal state weights are included. After a suffix
with initial state $i$ is fixed, its exact rewards are subtracted and
the prefix strip probability ending at $i$ is at most $C/\sqrt m$.
Multiplication by the suffix transition product and
$\pi_i^{-1}\le7/3$ gives the required conditional cylinder estimate.
Equivalently, detailed balance converts these products into the actual
stable-cylinder masses. The terminal coordinate belongs to its suffix
cylinder exactly, so this cover includes all coarea endpoints. The first
criterion applies again.

For the initial-age source use the second criterion with $q=(k,l)$,
$b_w=b_m$, $a_m=m^{-1/2}$, and the envelope
$C m^{-1}e^{-c|q-m\mu_*|^2/m}$ from
\eqref{eq:v60-cylinder-upper} at depth $L_m$.
At a coarea point set $\delta=t-3m-k-\epsilon l$. The implications are
\begin{align*}
 y_0<s&\ \Longrightarrow\ b_m\in[\delta,\delta+s],\\
 y_0>1-s&\ \Longrightarrow\
 b_m\in[\delta+1-s-\rho_0^m,\delta+1].
\end{align*}
Both intervals are retained. Summing their positive interval bounds
requires only the labels with $|\delta|\le1$. For each $l$ there
are at most three $k$, and the second-coordinate Gaussian sum is
$O(\sqrt m)$, exactly as in \eqref{eq:v60-strip-upper}.
This proves the age bound without summing a signed local-limit error.

Domination before pushforward proves the union and insertion statements.
Finally, \eqref{eq:v60-unweighted-pointwise} gives
$p_{m,\epsilon}(t)\ge c_M/\sqrt m$ on the stated central set for
large $m$, because the variances have common positive lower and upper
bounds. Divide the unchanged source density by this same-roof density.
Regular conditional probabilities are thereby specified for almost
every roof value. No roof set of positive measure is discarded.
\end{proof}

\subsection{A visible transition at the pressure scale}
The endpoint overlap is continuous and one-periodic. More precisely,
\begin{equation}\label{eq:v60-overlap-regularity}
 \|\Theta_{f,d}\|_\infty+
       \operatorname{Lip}_{\R/\Z}(\Theta_{f,d})
                        \le C\|f\|_{\rm Lip}\|d\|_{\rm Lip}.
\end{equation}
To verify this, extend $f$ periodically, allowing its one endpoint
jump. This extension has bounded variation on the circle. The overlap
is its circular correlation with the bounded periodic extension of
$d$. Translation in $L^1$ of a function of bounded variation is bounded
by the shift times its variation. This proves the Lipschitz estimate;
no equality of the values at $0,1$ is assumed.

Set
\begin{equation}\label{eq:v60-critical-variables}
 c_m=\epsilon\sqrt m,\qquad
 z=\frac{t-m\bar\tau_\epsilon}{\sqrt m},\qquad
 r_\epsilon=t-3m-\frac{2m\epsilon}{7},\qquad
 \beta_*=\frac{16}{17},\quad v_*=\frac{6}{119}.
\end{equation}
The last two constants are the conditional mean coefficient and variance
of the second coordinate of the centered $\Sigma_*$ Gaussian given its
first coordinate.

\begin{theorem}[The moving arithmetic profile on its critical scale]
\label{thm:v60-critical-profile}
For every $C_0<\infty$, uniformly for $|c_m|\le C_0$,
$|\epsilon|\le\epsilon_0$ and almost every $t\in\R$,
\begin{align}
 \sqrt m\,p_{m,\epsilon}^{f,d}(t)
 ={}&g_{204/343}(z)
       \int_\R\Theta_{f,d}(r_\epsilon-\beta_*c_m z-c_m v)
                                g_{v_*}(v)\,dv\notag\\
 &+O_{C_0}\left(\|f\|_{\rm Lip}\|d\|_{\rm Lip}
                      \frac{\log(2+m)^{3/2}}{\sqrt m}\right).
                                                  \label{eq:v60-critical-profile}
\end{align}
For the fixed physical witness functions
$f(y)=e^{2\pi iy}$ and $d(y)=e^{-2\pi iy}$, the main term is
\begin{equation}\label{eq:v60-witness-density}
 g_{204/343}(z)
 e^{-2\pi i(r_\epsilon-(16/17)c_m z)}
                      e^{-12\pi^2c_m^2/119}.
\end{equation}
Thus the moving damped peak and the pointwise density are evaluated on
the same nontrivial parameter family, together with the uniform positive
height estimates of Theorem~\ref{thm:v60-perturbed-height}.
\end{theorem}
\begin{proof}
Write the sum in \eqref{eq:v60-arithmetic-profile} on the grid
$y_l=(l-2m/7)/\sqrt m$. Its summand is
\[
 g_{\Sigma_*}(z-\epsilon y_l,y_l)
                           \Theta_{f,d}(r_\epsilon-c_m y_l).
\]
By \eqref{eq:v60-overlap-regularity} and Gaussian differentiation, its
variation as a function of $y$ is uniformly bounded by
$C_{C_0}\|f\|_{\rm Lip}\|d\|_{\rm Lip}$, independently of $z$ and
$r_\epsilon$. The Riemann error is therefore $O_{C_0}(m^{-1/2})$
with the same norm factor. Replacing $z-\epsilon y$ by $z$ in the
integral costs $O(|\epsilon|)$ uniformly in $z$: the integral of
$|y|\,|\partial_1g_{\Sigma_*}(z-\theta\epsilon y,y)|$ is uniformly
bounded for $0\le\theta\le1$. Since $|\epsilon|\le C_0/\sqrt m$,
this is another allowed error. Gaussian conditioning now gives
\[
 g_{\Sigma_*}(z,y)=g_{204/343}(z)
                                  g_{6/119}(y-(16/17)z).
\]
Combine this identity with \eqref{eq:v60-weighted-pointwise} to obtain
\eqref{eq:v60-critical-profile}.

For the displayed witness pair the original overlap equals
$\Theta_{f,d}(r)=e^{-2\pi ir}$. Its Gaussian integral is
$e^{-2\pi i(r_\epsilon-\beta_*c_mz)}e^{-2\pi^2v_*c_m^2}$,
which is \eqref{eq:v60-witness-density}. The damping coefficient is
exactly the one in \eqref{eq:v59-markov-expansions}:
$m\kappa_\epsilon=(12\pi^2/119)c_m^2+O_{C_0}(m^{-1/2})$.
This agreement is obtained from the original coarea density, not by
assuming that a spectral expansion already supplies a pointwise inverse.
\end{proof}

\paragraph{Relation to the original raw inversion problem.}
The scalar roof projection may change from lattice to nonlattice as
$\epsilon$ varies. The proof avoids any nonuniform scalar arithmetic
estimate by retaining the full integer pair until the positive coarea
sum is performed. The source-height criterion isolates the two inputs
which pressure alone cannot supply: a wordwise derivative and positive
local bounds after conditioning on the relevant cylinders. Finite-state
Markov additive local-limit methods are classical \cite{HL,HLuniform};
uniformity over compact transition families is also available there.
The specific additions here are the joint-label projection across an
arithmetic transition, its exact deterministic endpoint overlap, and
the three original positive source heights in that same family.

This verifies the new criterion in a correlated perturbed realization,
not in several unrelated non-Markov systems. In the triangular Lorentz
problem the labels $n,k,m$ and the next-collision clearance convention
must still be retained when verifying the criterion. Neither the stable
derivative in \eqref{eq:v60-exact-perturbed-coarea} nor the two-count
positive cylinder estimate is a theorem about those physical sources.
The two Lorentz heights in Corollary~\ref{cor:v51-positive-criterion},
and the unrestricted same-roof and essential-likelihood consequences
which require them, are not asserted here. All earlier exact-return
statements and their arithmetic zero classes remain in force.
